\documentclass[11pt]{article}

\usepackage[a4paper,margin=1in]{geometry}
\usepackage{amsmath,amssymb,amsthm,mathtools}
\usepackage[T1]{fontenc}
\usepackage{lmodern}
\usepackage{microtype}

% Fallback when hyperref is not loaded.
\providecommand{\texorpdfstring}[2]{#1}

\newtheorem{theorem}{Theorem}[section]
\newtheorem{lemma}[theorem]{Lemma}
\newtheorem{proposition}[theorem]{Proposition}
\newtheorem{corollary}[theorem]{Corollary}
\newtheorem{remark}[theorem]{Remark}
\newtheorem{conjecture}[theorem]{Conjecture}
\newtheorem{definition}[theorem]{Definition}

\newcommand{\C}{\mathbb C}
\newcommand{\R}{\mathbb R}
\newcommand{\HH}{\mathbb H}
\newcommand{\eps}{\varepsilon}

\title{A Dual-Tail Local Reduction for the Real-Slice Problem\\
in the Clean OBC Hatano--Nelson Family}
\author{}
\date{}

\begin{document}
\maketitle

\section{Setup}

Fix \(0<a<b\), and let
\[
A_n=\operatorname{tridiag}(a,0,b)\in\R^{n\times n}.
\]
For \(z=x+iy\in\C\), define
\[
B(z)\coloneqq zI-A_n,
\qquad
M_n(x,y)\coloneqq B(z)^*B(z),
\qquad
\lambda_*(x,y)\coloneqq \lambda_{\min}(M_n(x,y)).
\]
At a boundary point of the pseudospectrum,
\[
\Lambda=\lambda_*(x,y),
\]
one has
\[
M_n(x,y)-\Lambda I_n \succeq 0,
\qquad
\det(M_n(x,y)-\Lambda I_n)=0.
\]

We write
\[
q\coloneqq b-a>0,
\qquad
\eta\coloneqq ab,
\qquad
c\coloneqq (a+b)x-iqy,
\]
and also
\[
\delta\coloneqq x^2+y^2+a^2+b^2-\Lambda,
\qquad
\delta_a\coloneqq x^2+y^2+a^2-\Lambda=\delta-b^2,
\qquad
\delta_b\coloneqq x^2+y^2+b^2-\Lambda=\delta-a^2.
\]

A boundary point \((x,y)\) is called \emph{simple} if \(\Lambda\) is a simple eigenvalue of
\(M_n(x,y)\), equivalently
\[
\dim\ker\bigl(M_n(x,y)-\Lambda I_n\bigr)=1.
\]

The connectedness-equivalence problem is to prove that for every fixed \(x\in\R\) and every
\(\eps>0\), the vertical section
\[
\{y\in\R:\ \lambda_*(x,y)<\eps^2\}
\]
is either empty or an interval containing \(0\). At the simple-boundary level, it suffices to
prove that for every simple boundary point with \(y>0\),
\[
\partial_y \lambda_*(x,y)>0.
\]

\section{Diagonal gauge and weighted metric}

The non-selfadjoint asymmetry can be gauged away at the level of the hopping amplitudes, and the
corresponding weighted metric will reappear in the exact scalar identities below.

\begin{proposition}[Diagonal gauge, symmetrization, and weighted metric]\label{prop:gauge}
Set
\[
\rho\coloneqq \sqrt{\frac ab}\in(0,1),
\qquad
G\coloneqq \operatorname{diag}(1,\rho,\rho^2,\dots,\rho^{n-1}),
\]
and
\[
W\coloneqq G^{-2}
=
\operatorname{diag}(\omega_1,\dots,\omega_n),
\qquad
\omega_j\coloneqq \Bigl(\frac ba\Bigr)^{j-1}.
\]
Then
\[
G^{-1}A_nG=\sqrt{\eta}\,\operatorname{tridiag}(1,0,1),
\]
and
\[
WA_n=A_n^T W.
\]
Moreover,
\[
\omega_j\,\omega_{n+1-j}=\Bigl(\frac ba\Bigr)^{n-1}
\qquad (1\le j\le n).
\]
Consequently, if \(u=Gv\), then
\[
-\Im(u_j u_{n+1-j})
=
\rho^{\,n-1}\bigl(-\Im(v_j v_{n+1-j})\bigr),
\]
so the signs of the mirror-pair quantities are invariant under this diagonal gauge.
\end{proposition}

\begin{proof}
The superdiagonal entries of \(G^{-1}A_nG\) are
\[
\rho^{-(j-1)}\,b\,\rho^{j}=b\rho=\sqrt{ab}=\sqrt{\eta},
\]
and the subdiagonal entries are
\[
\rho^{-(j-1)}\,a\,\rho^{j-2}=a\rho^{-1}=\sqrt{ab}=\sqrt{\eta}.
\]
This proves the symmetrization formula.

For the weighted metric, note that
\[
b\,\omega_j=a\,\omega_{j+1},
\]
so the \((j,j+1)\)-entry of \(WA_n\) equals the \((j,j+1)\)-entry of \(A_n^T W\), and similarly
for the \((j+1,j)\)-entry. Thus \(WA_n=A_n^T W\).

Finally,
\[
\omega_j\omega_{n+1-j}
=
\Bigl(\frac ba\Bigr)^{j-1}\Bigl(\frac ba\Bigr)^{n-j}
=
\Bigl(\frac ba\Bigr)^{n-1}.
\]
If \(u=Gv\), then
\[
u_j u_{n+1-j}
=
\rho^{j-1}\rho^{n-j} v_j v_{n+1-j}
=
\rho^{n-1} v_j v_{n+1-j},
\]
and the sign statement follows.
\end{proof}

\section{Block form and dual tails}

We record the even case explicitly. The odd case has the same constant-coefficient interior and a
different terminal seed; all global identities in Sections~\ref{sec:takagi}--\ref{sec:status}
remain valid for arbitrary \(n\).

\subsection*{Even case \(n=2m\)}

Set
\[
w_k\coloneqq \binom{u_{2k-1}}{u_{2k}},
\qquad
1\le k\le m.
\]
Then
\[
\bigl(M_{2m}(x,y)-\Lambda I_{2m}\bigr)u=0
\]
is equivalent to the \(2\times2\) block Jacobi system
\[
B_L w_1 + A w_2 = 0,
\]
\[
A^* w_{k-1}+Bw_k+Aw_{k+1}=0,
\qquad
2\le k\le m-1,
\]
\[
A^*w_{m-1}+B_R w_m=0,
\]
where
\[
A=
\begin{pmatrix}
\eta & 0\\
-c & \eta
\end{pmatrix},
\qquad
\det A=\eta^2>0,
\]
and
\[
B_L=
\begin{pmatrix}
\delta_a & -c\\
-\bar c & \delta
\end{pmatrix},
\qquad
B=
\begin{pmatrix}
\delta & -c\\
-\bar c & \delta
\end{pmatrix},
\qquad
B_R=
\begin{pmatrix}
\delta & -c\\
-\bar c & \delta_b
\end{pmatrix}.
\]

\begin{lemma}[Proper leading and trailing subchains are positive definite]\label{lem:proper}
Let
\[
H\coloneqq M_{2m}(x,y)-\Lambda I_{2m}\succeq 0.
\]
For each \(2\le k\le m\), the trailing block principal submatrix of \(H\) on blocks
\(k,k+1,\dots,m\) is positive definite. For each \(1\le k\le m-1\), the leading block principal
submatrix of \(H\) on blocks \(1,2,\dots,k\) is positive definite.
\end{lemma}

\begin{proof}
We prove the trailing statement; the leading statement is the same with left and right reversed.
Let \(T_k\) be the trailing block principal submatrix on blocks \(k,\dots,m\), and suppose
\(T_k v=0\) for some nonzero block vector
\[
v=(\xi_k,\xi_{k+1},\dots,\xi_m)^T.
\]
Extend \(v\) by zeros to a block vector
\[
\widetilde v=(0,\dots,0,\xi_k,\xi_{k+1},\dots,\xi_m)^T
\]
in the full chain. Then
\[
\widetilde v^* H \widetilde v=v^*T_k v=0.
\]
Since \(H\succeq 0\), this implies \(H\widetilde v=0\). Looking at the \((k-1)\)-st block
equation of \(H\widetilde v=0\), we get
\[
A\xi_k=0.
\]
Because \(A\) is invertible, \(\xi_k=0\). Then the \(k\)-th block equation yields
\[
A\xi_{k+1}=0,
\]
hence \(\xi_{k+1}=0\). Continuing inductively gives \(\xi_j=0\) for all \(j\), contradicting
\(v\neq 0\). Therefore \(T_k\succ 0\).

The leading statement is identical, using the \((k+1)\)-st block equation and the invertibility
of \(A^*\).
\end{proof}

\begin{definition}[Right and left tail Weyl matrices]
In the even case, define recursively the right tails
\[
\Sigma_m\coloneqq B_R,
\qquad
\Sigma_k\coloneqq B_k-A\Sigma_{k+1}^{-1}A^*,
\qquad
k=m-1,\dots,1,
\]
where
\[
B_1\coloneqq B_L,
\qquad
B_k\coloneqq B
\quad (2\le k\le m-1),
\]
and the left tails
\[
\Gamma_1\coloneqq B_L,
\qquad
\Gamma_k\coloneqq B_k-A^*\Gamma_{k-1}^{-1}A,
\qquad
k=2,\dots,m,
\]
where
\[
B_k\coloneqq B
\quad (2\le k\le m-1),
\qquad
B_m\coloneqq B_R.
\]
\end{definition}

By Lemma~\ref{lem:proper}, all inverses appearing in the recursions are legitimate, and
\[
\Sigma_k\succ 0 \quad (2\le k\le m),
\qquad
\Gamma_k\succ 0 \quad (1\le k\le m-1).
\]
The endpoint tails \(\Sigma_1\) and \(\Gamma_m\) may be singular at the boundary.

\begin{proposition}[Tail action on the block nullchain]\label{prop:tails}
Let \(w=(w_1,\dots,w_m)\) be a nonzero block solution of the boundary equation
\[
\bigl(M_{2m}(x,y)-\Lambda I_{2m}\bigr)u=0.
\]
Then
\[
\Sigma_k w_k=-A^*w_{k-1},
\qquad
\Gamma_k w_k=-Aw_{k+1},
\]
for every admissible \(k\), with the endpoint conventions
\[
A^*w_0=0,
\qquad
Aw_{m+1}=0.
\]
\end{proposition}

\begin{proof}
For the right tails, the identity at \(k=m\) is exactly the terminal equation
\[
\Sigma_m w_m=B_R w_m=-A^*w_{m-1}.
\]
Assume \(\Sigma_{k+1}w_{k+1}=-A^*w_k\). Since \(\Sigma_{k+1}\) is invertible for \(k+1\ge 2\),
\[
w_{k+1}=-\Sigma_{k+1}^{-1}A^*w_k.
\]
Substituting into the \(k\)-th block equation gives
\[
A^*w_{k-1}+\bigl(B_k-A\Sigma_{k+1}^{-1}A^*\bigr)w_k=0,
\]
that is,
\[
\Sigma_k w_k=-A^*w_{k-1}.
\]
Backward induction proves the right-tail identity for all \(k\). The left-tail identity is
proved similarly, starting from
\[
\Gamma_1 w_1=B_L w_1=-Aw_2.
\]
\end{proof}

\begin{lemma}[No block of the nullchain vanishes]\label{lem:blocknonzero}
If \(w=(w_1,\dots,w_m)\) is a nonzero block solution of the boundary equation, then
\[
w_k\neq 0
\qquad
(1\le k\le m).
\]
\end{lemma}

\begin{proof}
If \(w_k=0\), then Proposition~\ref{prop:tails} gives
\[
A^*w_{k-1}=0,
\qquad
Aw_{k+1}=0,
\]
hence
\[
w_{k-1}=0,
\qquad
w_{k+1}=0,
\]
because \(A\) is invertible. Iterating in both directions yields \(w_j=0\) for all \(j\), a
contradiction.
\end{proof}

\begin{definition}[Local effective block]
For \(1\le k\le m\), define
\[
K_k\coloneqq \Sigma_k+\Gamma_k-B_k,
\]
where \(B_1=B_L\), \(B_m=B_R\), and \(B_k=B\) for \(2\le k\le m-1\).
\end{definition}

\begin{proposition}[Exact local invariant]\label{prop:K}
At a simple boundary point \(\Lambda=\lambda_*(x,y)\), one has
\[
K_k\succeq 0,
\qquad
\operatorname{rank}K_k=1,
\qquad
K_k w_k=0
\]
for every \(k\). Therefore
\[
K_k=
\begin{pmatrix}
A_k & \zeta_k+i t_k\\
\zeta_k-i t_k & D_k
\end{pmatrix},
\]
with
\[
A_k\ge 0,\qquad D_k\ge 0,\qquad \zeta_k,t_k\in\R,
\qquad
A_kD_k=\zeta_k^2+t_k^2.
\]
\end{proposition}

\begin{proof}
Eliminate all blocks except block \(k\) from the positive semidefinite matrix
\[
H=M_{2m}(x,y)-\Lambda I_{2m}
\]
by Schur complement. Lemma~\ref{lem:proper} shows that the complementary leading/trailing block
submatrices are positive definite, so this Schur complement is legitimate and equals \(K_k\).
Hence \(K_k\succeq 0\).

Since \(H\) has one-dimensional kernel and \(w_k\neq 0\) by Lemma~\ref{lem:blocknonzero}, the
kernel of the local Schur complement is exactly \(\C w_k\). Thus \(\dim\ker K_k=1\), so
\(\operatorname{rank}K_k=1\), and \(K_k w_k=0\).

The determinant identity
\[
A_kD_k=\zeta_k^2+t_k^2
\]
is simply \(\det K_k=0\).
\end{proof}

\section{Tail imaginary parts and orientation of the local current}

Write
\[
\Sigma_k=
\begin{pmatrix}
u_k^R & \beta_k^R\\
\overline{\beta_k^R} & v_k^R
\end{pmatrix},
\qquad
\beta_k^R=b_k^R+i\tau_k^R,
\qquad
\Delta_k^R\coloneqq u_k^R v_k^R-\bigl|\beta_k^R\bigr|^2,
\]
and
\[
\Gamma_k=
\begin{pmatrix}
U_k^L & \gamma_k^L\\
\overline{\gamma_k^L} & V_k^L
\end{pmatrix},
\qquad
\gamma_k^L=g_k^L+i\tau_k^L,
\qquad
\Delta_k^L\coloneqq U_k^L V_k^L-\bigl|\gamma_k^L\bigr|^2.
\]

\begin{proposition}[Right and left tail recurrences]\label{prop:taurec}
For \(1\le k\le m-1\), one has
\[
u_k^R
=
\delta-\frac{\eta^2 v_{k+1}^R}{\Delta_{k+1}^R},
\]
and
\[
\beta_k^R
=
-c+\frac{\eta(\eta\beta_{k+1}^R+v_{k+1}^R\bar c)}{\Delta_{k+1}^R}.
\]
In particular,
\[
\tau_k^R
=
qy+\frac{\eta(\eta\tau_{k+1}^R+qy\,v_{k+1}^R)}{\Delta_{k+1}^R}.
\]
Similarly, for \(2\le k\le m\), one has
\[
V_k^L
=
\delta-\frac{\eta^2 U_{k-1}^L}{\Delta_{k-1}^L},
\]
and
\[
\gamma_k^L
=
-c+\frac{\eta(\eta\gamma_{k-1}^L+U_{k-1}^L\bar c)}{\Delta_{k-1}^L}.
\]
In particular,
\[
\tau_k^L
=
qy+\frac{\eta(\eta\tau_{k-1}^L+qy\,U_{k-1}^L)}{\Delta_{k-1}^L}.
\]
\end{proposition}

\begin{proof}
For the right tails,
\[
\Sigma_k=B_k-A\Sigma_{k+1}^{-1}A^*,
\qquad
\Sigma_{k+1}^{-1}
=
\frac1{\Delta_{k+1}^R}
\begin{pmatrix}
v_{k+1}^R & -\beta_{k+1}^R\\
-\overline{\beta_{k+1}^R} & u_{k+1}^R
\end{pmatrix}.
\]
A direct multiplication gives
\[
A\Sigma_{k+1}^{-1}A^*
=
\frac1{\Delta_{k+1}^R}
\begin{pmatrix}
\eta^2 v_{k+1}^R &
-\eta\bigl(\eta\beta_{k+1}^R+v_{k+1}^R\bar c\bigr)
\\[1ex]
-\eta\bigl(\eta\overline{\beta_{k+1}^R}+v_{k+1}^R c\bigr) &
*
\end{pmatrix}.
\]
Subtracting from \(B_k\) yields the claimed formulas for \(u_k^R\) and \(\beta_k^R\), and
taking imaginary parts gives the recurrence for \(\tau_k^R\).

The left-tail formulas are identical, starting from
\[
\Gamma_k=B_k-A^*\Gamma_{k-1}^{-1}A,
\qquad
\Gamma_{k-1}^{-1}
=
\frac1{\Delta_{k-1}^L}
\begin{pmatrix}
V_{k-1}^L & -\gamma_{k-1}^L\\
-\overline{\gamma_{k-1}^L} & U_{k-1}^L
\end{pmatrix}.
\]
One obtains
\[
A^*\Gamma_{k-1}^{-1}A
=
\frac1{\Delta_{k-1}^L}
\begin{pmatrix}
* &
-\eta\bigl(\eta\gamma_{k-1}^L+U_{k-1}^L\bar c\bigr)
\\[1ex]
-\eta\bigl(\eta\overline{\gamma_{k-1}^L}+U_{k-1}^L c\bigr) &
\eta^2 U_{k-1}^L
\end{pmatrix},
\]
and the conclusions follow.
\end{proof}

\begin{corollary}[Strict positivity of the tail imaginary parts]\label{cor:taupos}
One has
\[
\tau_m^R=qy,
\qquad
\tau_k^R>qy
\quad (1\le k\le m-1),
\]
and
\[
\tau_1^L=qy,
\qquad
\tau_k^L>qy
\quad (2\le k\le m).
\]
\end{corollary}

\begin{proof}
The seeds are
\[
\Sigma_m=B_R,
\qquad
\Gamma_1=B_L,
\]
so
\[
\tau_m^R=\tau_1^L=\Im(-c)=qy.
\]
For \(k\le m-1\), Proposition~\ref{prop:taurec} gives
\[
\tau_k^R
=
qy+\frac{\eta(\eta\tau_{k+1}^R+qy\,v_{k+1}^R)}{\Delta_{k+1}^R}.
\]
Since \(\Sigma_{k+1}\succ 0\), one has
\[
v_{k+1}^R>0,
\qquad
\Delta_{k+1}^R>0,
\]
and the extra term is strictly positive. Hence \(\tau_k^R>qy\). The left-tail statement is
proved identically using \(\Gamma_{k-1}\succ 0\).
\end{proof}

\begin{proposition}[Orientation of the local effective block]\label{prop:orientation}
For every \(1\le k\le m\), the off-diagonal imaginary part of \(K_k\) is
\[
t_k=\tau_k^R+\tau_k^L-qy,
\]
and in particular
\[
t_k>0.
\]
Consequently,
\[
A_k>0,
\qquad
D_k>0.
\]
Moreover, if
\[
r_k\coloneqq \frac{w_{k,2}}{w_{k,1}},
\]
then \(r_k\) is well defined and
\[
r_k\in\HH:=\{z\in\C:\Im z>0\}.
\]
Finally, the within-block current
\[
j_k\coloneqq \Im(\overline{w_{k,1}}\,w_{k,2})
\]
satisfies
\[
j_k=\frac{t_k}{D_k}|w_{k,1}|^2=\frac{t_k}{A_k}|w_{k,2}|^2>0.
\]
\end{proposition}

\begin{proof}
The off-diagonal entry of every \(B_k\) is \(-c\), whose imaginary part is \(qy\). Hence
\[
t_k=\Im\bigl((K_k)_{12}\bigr)=\tau_k^R+\tau_k^L-qy.
\]
By Corollary~\ref{cor:taupos}, at least one of the two terms \(\tau_k^R,\tau_k^L\) is strictly
larger than \(qy\), while the other is at least \(qy\). Therefore \(t_k>0\).

Since \(K_k\succeq 0\), \(\operatorname{rank}K_k=1\), and its off-diagonal entry has nonzero
imaginary part, both diagonal entries must be strictly positive:
\[
A_k>0,\qquad D_k>0.
\]
Indeed, if one diagonal entry vanished, positivity would force the entire corresponding row and
column to vanish, contradicting \(t_k>0\).

Now \(K_k w_k=0\) and \(t_k>0\) imply that neither component of \(w_k\) can vanish. Thus
\(r_k\) is well defined. From the second row of \(K_k w_k=0\),
\[
(\zeta_k-i t_k)w_{k,1}+D_k w_{k,2}=0,
\]
so
\[
r_k=-\frac{\zeta_k-i t_k}{D_k},
\qquad
\Im r_k=\frac{t_k}{D_k}>0.
\]
This proves \(r_k\in\HH\). Finally,
\[
j_k=\Im(\overline{w_{k,1}}\,w_{k,2})
=\Im(r_k)|w_{k,1}|^2
=\frac{t_k}{D_k}|w_{k,1}|^2>0.
\]
Using the first row of \(K_k w_k=0\),
\[
A_k w_{k,1}+(\zeta_k+i t_k)w_{k,2}=0,
\]
one similarly gets
\[
\frac{1}{r_k}=-\frac{\zeta_k+i t_k}{A_k},
\qquad
j_k=\frac{t_k}{A_k}|w_{k,2}|^2.
\]
\end{proof}

\section{Cross-block currents and local mirror criteria}

Define the cross-block currents
\[
h_k\coloneqq \Im(\overline{w_{k,2}}\,w_{k+1,1}),
\qquad
1\le k\le m-1.
\]

\begin{proposition}[Exact formulas for the adjacent cross-currents]\label{prop:h}
For every \(1\le k\le m-1\), one has
\[
h_k
=
\frac{
\eta v_{k+1}^R\,j_k
+
(\eta\tau_{k+1}^R+qy\,v_{k+1}^R)\,|w_{k,2}|^2
}{
\Delta_{k+1}^R
}.
\]
For every \(2\le k\le m\), one has
\[
h_{k-1}
=
\frac{
\eta U_{k-1}^L\,j_k
+
(\eta\tau_{k-1}^L+qy\,U_{k-1}^L)\,|w_{k,1}|^2
}{
\Delta_{k-1}^L
}.
\]
In particular,
\[
h_k>0
\qquad
(1\le k\le m-1).
\]
\end{proposition}

\begin{proof}
From the right-tail identity
\[
w_{k+1}=-\Sigma_{k+1}^{-1}A^*w_k,
\]
with
\[
\Sigma_{k+1}^{-1}
=
\frac1{\Delta_{k+1}^R}
\begin{pmatrix}
v_{k+1}^R & -\beta_{k+1}^R\\
-\overline{\beta_{k+1}^R} & u_{k+1}^R
\end{pmatrix},
\qquad
A^*=
\begin{pmatrix}
\eta & -\bar c\\
0 & \eta
\end{pmatrix},
\]
we obtain
\[
w_{k+1,1}
=
-\frac{
\eta v_{k+1}^R\,w_{k,1}
-(v_{k+1}^R\bar c+\eta\beta_{k+1}^R)w_{k,2}
}{
\Delta_{k+1}^R
}.
\]
Therefore
\[
\overline{w_{k,2}}\,w_{k+1,1}
=
-\frac{
\eta v_{k+1}^R\,\overline{w_{k,2}}\,w_{k,1}
-(v_{k+1}^R\bar c+\eta\beta_{k+1}^R)|w_{k,2}|^2
}{
\Delta_{k+1}^R
}.
\]
Taking imaginary parts and using
\[
\Im(\overline{w_{k,2}}\,w_{k,1})=-j_k,
\qquad
\Im(\bar c)=qy,
\qquad
\Im(\beta_{k+1}^R)=\tau_{k+1}^R,
\]
yields
\[
h_k
=
\frac{
\eta v_{k+1}^R\,j_k
+
(\eta\tau_{k+1}^R+qy\,v_{k+1}^R)|w_{k,2}|^2
}{
\Delta_{k+1}^R
}.
\]

Similarly, from the left-tail identity
\[
w_{k-1}=-(\Gamma_{k-1}^{-1}A)w_k,
\]
with
\[
\Gamma_{k-1}^{-1}
=
\frac1{\Delta_{k-1}^L}
\begin{pmatrix}
V_{k-1}^L & -\gamma_{k-1}^L\\
-\overline{\gamma_{k-1}^L} & U_{k-1}^L
\end{pmatrix},
\qquad
A=
\begin{pmatrix}
\eta & 0\\
-c & \eta
\end{pmatrix},
\]
we obtain
\[
w_{k-1,2}
=
-\frac{
(-\eta\overline{\gamma_{k-1}^L}-U_{k-1}^L c)\,w_{k,1}
+\eta U_{k-1}^L\,w_{k,2}
}{
\Delta_{k-1}^L
}.
\]
Hence
\[
\overline{w_{k-1,2}}\,w_{k,1}
=
\frac{
(\eta\gamma_{k-1}^L+U_{k-1}^L\bar c)|w_{k,1}|^2
-\eta U_{k-1}^L\,\overline{w_{k,2}}\,w_{k,1}
}{
\Delta_{k-1}^L
}.
\]
Taking imaginary parts gives
\[
h_{k-1}
=
\frac{
(\eta\tau_{k-1}^L+qy\,U_{k-1}^L)|w_{k,1}|^2
+\eta U_{k-1}^L\,j_k
}{
\Delta_{k-1}^L
}.
\]

Finally, the positivity follows from Proposition~\ref{prop:orientation},
Corollary~\ref{cor:taupos}, and the positive definiteness of the tails:
\[
j_k>0,\qquad
v_{k+1}^R>0,\qquad
U_{k-1}^L>0,\qquad
\Delta_{k+1}^R>0,\qquad
\Delta_{k-1}^L>0.
\]
\end{proof}

\begin{corollary}[Positivity of all site currents in the even case]\label{cor:sitecurr}
Define the site currents
\[
I_{2k-1}\coloneqq j_k
\qquad (1\le k\le m),
\qquad
I_{2k}\coloneqq h_k
\qquad (1\le k\le m-1).
\]
Then
\[
I_j>0
\qquad
(1\le j\le 2m-1).
\]
Consequently,
\[
s\,p_{2m}=y|u_{2m}|^2+aI_{2m-1}>0,
\]
and therefore
\[
p_1=p_{2m}>0.
\]
\end{corollary}

\begin{proof}
The positivity of \(j_k\) is Proposition~\ref{prop:orientation}, and the positivity of \(h_k\) is
Proposition~\ref{prop:h}. The endpoint identity is the \(j=2m\) case of the Takagi site
identity from Proposition~\ref{prop:Takagi} below:
\[
s\,p_{2m}=y|u_{2m}|^2+aI_{2m-1}-bI_{2m},
\qquad
I_{2m}=0.
\]
\end{proof}

\begin{proposition}[Local criteria for mirror-pair positivity]\label{prop:localcriteria}
For every interior block \(2\le k\le m-1\), define
\[
p_j\coloneqq -\Im(u_j u_{2m+1-j}),
\qquad
s\coloneqq \sigma_{\min}(x+iy-A_{2m}).
\]
Then
\[
s\,p_{2k-1}
=
y|w_{k,1}|^2+a h_{k-1}-b j_k,
\]
\[
s\,p_{2k}
=
y|w_{k,2}|^2+a j_k-b h_k.
\]
Consequently,
\[
s\,p_{2k-1}=|w_{k,1}|^2\,\Psi_k^-,
\qquad
s\,p_{2k}=|w_{k,2}|^2\,\Psi_k^+,
\]
where
\[
\Psi_k^-
=
y
+
\frac{a(\eta\tau_{k-1}^L+qy\,U_{k-1}^L)}{\Delta_{k-1}^L}
+
\left(
\frac{a\eta U_{k-1}^L}{\Delta_{k-1}^L}-b
\right)\frac{t_k}{D_k},
\]
and
\[
\Psi_k^+
=
y
-
\frac{b(\eta\tau_{k+1}^R+qy\,v_{k+1}^R)}{\Delta_{k+1}^R}
+
\left(
a-\frac{b\eta v_{k+1}^R}{\Delta_{k+1}^R}
\right)\frac{t_k}{A_k}.
\]
Equivalently,
\[
p_{2k-1}>0 \iff \Psi_k^- >0,
\qquad
p_{2k}>0 \iff \Psi_k^+ >0.
\]
\end{proposition}

\begin{proof}
The Takagi site identities from Proposition~\ref{prop:Takagi} give
\[
s\,p_j
=
y|u_j|^2+a I_{j-1}-b I_j.
\]
For \(j=2k-1\), this becomes
\[
s\,p_{2k-1}
=
y|w_{k,1}|^2+a h_{k-1}-b j_k.
\]
For \(j=2k\), it becomes
\[
s\,p_{2k}
=
y|w_{k,2}|^2+a j_k-b h_k.
\]
Now substitute the formulas for \(h_{k-1}\) and \(h_k\) from Proposition~\ref{prop:h}, and use
Proposition~\ref{prop:orientation}:
\[
j_k=\frac{t_k}{D_k}|w_{k,1}|^2=\frac{t_k}{A_k}|w_{k,2}|^2.
\]
Factoring out \(|w_{k,1}|^2\) and \(|w_{k,2}|^2\) yields the displayed formulas for
\(\Psi_k^\pm\).
\end{proof}

\begin{proposition}[Same-level rewrite of \(\Psi_k^\pm\)]\label{prop:samelevel}
For every interior block \(2\le k\le m-1\), one has
\[
aA_k\,\Psi_k^+
=
\bigl(u_k^R-\delta_b\bigr)\bigl(ay+\tau_k^L\bigr)
+
\bigl(by-\tau_k^R\bigr)\bigl(U_k^L-a^2\bigr),
\]
and
\[
bD_k\,\Psi_k^-
=
\bigl(by-\tau_k^R\bigr)\bigl(V_k^L-\delta_a\bigr)
+
\bigl(ay+\tau_k^L\bigr)\bigl(v_k^R-b^2\bigr).
\]
Equivalently,
\[
\Psi_k^+>0
\iff
\bigl(u_k^R-\delta_b\bigr)\bigl(ay+\tau_k^L\bigr)
+
\bigl(by-\tau_k^R\bigr)\bigl(U_k^L-a^2\bigr)>0,
\]
and
\[
\Psi_k^->0
\iff
\bigl(by-\tau_k^R\bigr)\bigl(V_k^L-\delta_a\bigr)
+
\bigl(ay+\tau_k^L\bigr)\bigl(v_k^R-b^2\bigr)>0.
\]
\end{proposition}

\begin{proof}
From Proposition~\ref{prop:taurec},
\[
u_k^R-\delta_b
=
a^2-\frac{\eta^2 v_{k+1}^R}{\Delta_{k+1}^R}
=
a\left(a-\frac{b\eta v_{k+1}^R}{\Delta_{k+1}^R}\right),
\]
and
\[
by-\tau_k^R
=
by-qy-\frac{\eta(\eta\tau_{k+1}^R+qy\,v_{k+1}^R)}{\Delta_{k+1}^R}
=
a\left(
y-\frac{b(\eta\tau_{k+1}^R+qy\,v_{k+1}^R)}{\Delta_{k+1}^R}
\right).
\]
Also,
\[
A_k
=
u_k^R+U_k^L-\delta
=
\bigl(u_k^R-\delta_b\bigr)+\bigl(U_k^L-a^2\bigr),
\]
and by Proposition~\ref{prop:orientation},
\[
t_k
=
\tau_k^R+\tau_k^L-qy
=
\bigl(ay+\tau_k^L\bigr)-\bigl(by-\tau_k^R\bigr).
\]
Hence
\[
a\Psi_k^+
=
\bigl(by-\tau_k^R\bigr)
+
\frac{u_k^R-\delta_b}{A_k}\,t_k,
\]
so
\[
aA_k\,\Psi_k^+
=
A_k\bigl(by-\tau_k^R\bigr)+\bigl(u_k^R-\delta_b\bigr)t_k.
\]
Substituting the expressions for \(A_k\) and \(t_k\) gives
\[
aA_k\,\Psi_k^+
=
\bigl(u_k^R-\delta_b\bigr)\bigl(ay+\tau_k^L\bigr)
+
\bigl(by-\tau_k^R\bigr)\bigl(U_k^L-a^2\bigr).
\]

Similarly, Proposition~\ref{prop:taurec} gives
\[
V_k^L-\delta_a
=
b^2-\frac{\eta^2 U_{k-1}^L}{\Delta_{k-1}^L}
=
b\left(b-\frac{a\eta U_{k-1}^L}{\Delta_{k-1}^L}\right),
\]
and
\[
ay+\tau_k^L
=
ay+qy+\frac{\eta(\eta\tau_{k-1}^L+qy\,U_{k-1}^L)}{\Delta_{k-1}^L}
=
b\left(
y+\frac{a(\eta\tau_{k-1}^L+qy\,U_{k-1}^L)}{\Delta_{k-1}^L}
\right).
\]
Moreover,
\[
D_k
=
v_k^R+V_k^L-\delta
=
\bigl(v_k^R-b^2\bigr)+\bigl(V_k^L-\delta_a\bigr),
\]
and
\[
t_k
=
\bigl(ay+\tau_k^L\bigr)-\bigl(by-\tau_k^R\bigr).
\]
Therefore
\[
b\Psi_k^-
=
\bigl(ay+\tau_k^L\bigr)
-
\frac{V_k^L-\delta_a}{D_k}\,t_k,
\]
hence
\[
bD_k\,\Psi_k^-
=
D_k\bigl(ay+\tau_k^L\bigr)-\bigl(V_k^L-\delta_a\bigr)t_k.
\]
Substituting the expressions for \(D_k\) and \(t_k\) yields
\[
bD_k\,\Psi_k^-
=
\bigl(by-\tau_k^R\bigr)\bigl(V_k^L-\delta_a\bigr)
+
\bigl(ay+\tau_k^L\bigr)\bigl(v_k^R-b^2\bigr).
\]
\end{proof}

\begin{remark}
The local mirror-pair sign problem has therefore been reduced to a same-level pairing of the
left and right tails. In this form, the quantities \(t_k/A_k\) and \(t_k/D_k\) disappear, and
the remaining gap becomes a cone-pairing problem for the dual Weyl data at a single block.
\end{remark}

\section{Global Takagi identities and the weighted exact invariant}\label{sec:takagi}

We now record the cleanest global scalar reductions.

\begin{proposition}[Takagi form and exact current identities]\label{prop:Takagi}
Let
\[
B(z)\coloneqq zI-A_n,
\qquad
z=x+iy,
\qquad
y>0.
\]
Choose a right singular vector \(u=(u_1,\dots,u_n)^T\neq 0\) in Takagi form:
\[
B(z)u=s\,J\overline{u},
\qquad
s=\sigma_{\min}(B(z)),
\]
where \(J\) is the reversal matrix.

Define
\[
I_j\coloneqq \Im(\overline{u_j}u_{j+1}),
\qquad
1\le j\le n-1,
\qquad
I_0=I_n=0,
\]
and
\[
p_j\coloneqq -\Im(u_j u_{n+1-j}),
\qquad
1\le j\le n.
\]
Then for every \(j=1,\dots,n\),
\[
s\,p_j
=
y|u_j|^2+a I_{j-1}-b I_j.
\]
If moreover \(\|u\|=1\), then
\[
s\sum_{j=1}^n p_j
=
y-q\sum_{j=1}^{n-1} I_j.
\]
Consequently,
\[
\partial_y \lambda_*(x,y)
=
2y-2q\sum_{j=1}^{n-1} I_j
=
2s\sum_{j=1}^n p_j.
\]
\end{proposition}

\begin{proof}
The Takagi equation is
\[
(x+iy)u_j-a u_{j-1}-b u_{j+1}=s\,\overline{u_{n+1-j}},
\qquad
u_0=u_{n+1}=0.
\]
Multiplying by \(\overline{u_j}\) and taking imaginary parts gives
\[
s\,p_j=y|u_j|^2+a I_{j-1}-b I_j.
\]
Summing over \(j\) gives
\[
s\sum_{j=1}^n p_j
=
y\sum_{j=1}^n |u_j|^2
+a\sum_{j=1}^n I_{j-1}
-b\sum_{j=1}^n I_j
=
y-q\sum_{j=1}^{n-1} I_j.
\]
For normalized \(u\), the Hellmann--Feynman formula yields
\[
\partial_y \lambda_*(x,y)=2y-2q\sum_{j=1}^{n-1} I_j,
\]
hence
\[
\partial_y \lambda_*(x,y)=2s\sum_{j=1}^n p_j.
\]
\end{proof}

\begin{proposition}[Weighted exact invariant]\label{prop:weighted}
Let
\[
\omega_j\coloneqq \Bigl(\frac ba\Bigr)^{j-1},
\qquad
1\le j\le n,
\]
so that \(\omega_j=W_{jj}\) are the diagonal entries of the weighted metric \(W\) from
Proposition~\ref{prop:gauge}. Then for every \(1\le \ell\le r\le n\),
\[
s\sum_{j=\ell}^r \omega_j p_j
=
y\sum_{j=\ell}^r \omega_j |u_j|^2
+
a\omega_\ell I_{\ell-1}
-
b\omega_r I_r.
\]
In particular,
\[
s\sum_{j=1}^n \omega_j p_j
=
y\sum_{j=1}^n \omega_j |u_j|^2
>0.
\]
\end{proposition}

\begin{proof}
Multiply the site identity from Proposition~\ref{prop:Takagi} by \(\omega_j\):
\[
s\,\omega_j p_j
=
y\omega_j |u_j|^2
+
a\omega_j I_{j-1}
-
b\omega_j I_j.
\]
Since
\[
a\omega_j=b\omega_{j-1}
\qquad (j\ge 2),
\]
the sum over \(j=\ell,\dots,r\) telescopes:
\[
a\sum_{j=\ell}^r \omega_j I_{j-1}
-
b\sum_{j=\ell}^r \omega_j I_j
=
a\omega_\ell I_{\ell-1}
-
b\omega_r I_r.
\]
This gives the first identity. Setting \(\ell=1\) and \(r=n\), and using \(I_0=I_n=0\), gives
\[
s\sum_{j=1}^n \omega_j p_j
=
y\sum_{j=1}^n \omega_j |u_j|^2,
\]
which is strictly positive because \(y>0\) and \(u\neq 0\).
\end{proof}

\begin{corollary}[Positive weighted suffix sums in the even case]\label{cor:weightsuffix}
Assume \(n=2m\). Then for every \(1\le \ell\le 2m\),
\[
s\sum_{j=\ell}^{2m} \omega_j p_j
=
y\sum_{j=\ell}^{2m}\omega_j |u_j|^2
+
a\omega_\ell I_{\ell-1}
>0.
\]
\end{corollary}

\begin{proof}
This is Proposition~\ref{prop:weighted} with \(r=2m\), together with
\[
I_{2m}=0
\]
and the positivity of all site currents \(I_1,\dots,I_{2m-1}\) from
Corollary~\ref{cor:sitecurr}.
\end{proof}

\section{Connectedness equivalence: exact reduction and remaining gap}

\begin{lemma}[Topological vertical-slice criterion]\label{lem:topological}
Let \(\Omega\subset\C\) be open. If every vertical slice
\[
\{y\in\R:\ x+iy\in\Omega\}
\]
is either empty or an interval containing \(0\), then
\[
\Omega\ \text{is connected}
\quad\Longleftrightarrow\quad
\Omega\cap\R\ \text{is connected}.
\]
\end{lemma}

\begin{conjecture}[Mirror-pair positivity]\label{conj:mirror}
At every simple boundary point \(\Lambda=\lambda_*(x,y)\) with \(y>0\), one has
\[
p_j>0
\qquad (1\le j\le n).
\]
\end{conjecture}

\begin{conjecture}[Even-\(n\) local tail-pair positivity]\label{conj:tailpair}
Assume \(n=2m\). At every simple boundary point with \(y>0\), for each interior block
\(2\le k\le m-1\),
\[
\bigl(u_k^R-\delta_b\bigr)\bigl(ay+\tau_k^L\bigr)
+
\bigl(by-\tau_k^R\bigr)\bigl(U_k^L-a^2\bigr)>0,
\]
and
\[
\bigl(by-\tau_k^R\bigr)\bigl(V_k^L-\delta_a\bigr)
+
\bigl(ay+\tau_k^L\bigr)\bigl(v_k^R-b^2\bigr)>0.
\]
\end{conjecture}

\begin{remark}
By Propositions~\ref{prop:localcriteria} and \ref{prop:samelevel}, the even-\(n\) conjecture is
equivalent to
\[
\Psi_k^- >0,
\qquad
\Psi_k^+ >0
\qquad
(2\le k\le m-1),
\]
and hence to positivity of the interior mirror pairs
\[
p_{2k-1},\quad p_{2k}.
\]
The endpoint pair is already positive in the even case by Corollary~\ref{cor:sitecurr}.
\end{remark}

\begin{theorem}[Conditional simple-boundary completion]\label{thm:conditional}
Assume Conjecture~\ref{conj:mirror}. Then at every simple boundary point with \(y>0\),
\[
\partial_y \lambda_*(x,y)>0.
\]
Consequently, if all positive boundary points on the level \(\lambda_*=\eps^2\) are simple
(for example, on generic levels), then every vertical slice of \(\sigma_\eps(A_n)\) is either
empty or an interval containing \(0\). In that case,
\[
\sigma_\eps(A_n)\ \text{is connected}
\quad\Longleftrightarrow\quad
\sigma_\eps(A_n)\cap\R\ \text{is connected}.
\]
\end{theorem}

\begin{proof}
Let \((x,y)\) be a simple boundary point with \(y>0\), and let \(u\) be a normalized Takagi
singular vector. Proposition~\ref{prop:Takagi} gives
\[
\partial_y \lambda_*(x,y)=2s\sum_{j=1}^n p_j.
\]
Under Conjecture~\ref{conj:mirror}, the sum on the right is strictly positive, so
\[
\partial_y \lambda_*(x,y)>0.
\]

Now fix \(x\in\R\) and define
\[
F(y)\coloneqq \lambda_*(x,y)-\eps^2.
\]
Then \(F\) is even, continuous, and \(F(y)\to+\infty\) as \(|y|\to\infty\). If every positive
boundary point \(F(y)=0\) is simple, then each positive root is crossed with positive derivative.
Hence
\[
\{y\ge 0:\ F(y)<0\}
\]
is either empty or an interval of the form \([0,y_*)\). By evenness,
\[
\{y\in\R:\ F(y)<0\}
\]
is either empty or an interval containing \(0\). Thus every vertical slice of
\(\sigma_\eps(A_n)\) has the required form, and Lemma~\ref{lem:topological} proves the
connectedness equivalence.
\end{proof}

\begin{remark}
The theorem only needs positivity of the unweighted mirror sum
\[
\sum_{j=1}^n p_j.
\]
The stronger pointwise conjecture is one sufficient route. By contrast, the weighted exact
invariant already proves
\[
\sum_{j=1}^n \omega_j p_j>0
\]
without further work. Thus an alternative route to connectedness would be to compare the
weighted and unweighted mirror sums, using the positivity of the site currents and the local
tail-pair structure.
\end{remark}

\section{Current status}\label{sec:status}

\begin{remark}
The exact progress is now the following.

\smallskip
\noindent
\textup{(1)} The non-selfadjoint hopping is controlled by the diagonal gauge
\[
G^{-1}A_nG=\sqrt{\eta}\,\operatorname{tridiag}(1,0,1),
\]
and by the weighted metric
\[
W=\operatorname{diag}\Bigl(1,\frac ba,\Bigl(\frac ba\Bigr)^2,\dots,
\Bigl(\frac ba\Bigr)^{n-1}\Bigr),
\qquad
WA_n=A_n^T W.
\]
The mirror-pair signs are invariant under this gauge.

\smallskip
\noindent
\textup{(2)} In the even block reduction, every proper leading or trailing principal block
submatrix is positive definite. Hence the parity tails \(\Sigma_k\) and \(\Gamma_k\) are
well-defined by genuine Schur complementation.

\smallskip
\noindent
\textup{(3)} The local Schur complement
\[
K_k=\Sigma_k+\Gamma_k-B_k
\]
is positive semidefinite of rank \(1\), annihilates the local block \(w_k\), and has
off-diagonal imaginary part
\[
t_k=\tau_k^R+\tau_k^L-qy.
\]
The exact tail recurrences imply
\[
t_k>0,
\]
so
\[
A_k>0,\qquad D_k>0,\qquad r_k\in\HH,\qquad j_k>0.
\]

\smallskip
\noindent
\textup{(4)} The adjacent cross-currents \(h_k\) are explicit affine functions of the local
current \(j_k\) and the neighboring tails, and are strictly positive. Thus in the even case all
site currents
\[
I_j=\Im(\overline{u_j}u_{j+1})
\]
are positive.

\smallskip
\noindent
\textup{(5)} The mirror-pair sign problem has been reduced to the local inequalities
\[
\Psi_k^- >0,\qquad \Psi_k^+>0,
\]
and then further rewritten as the same-level tail pairings
\[
\bigl(u_k^R-\delta_b\bigr)\bigl(ay+\tau_k^L\bigr)
+
\bigl(by-\tau_k^R\bigr)\bigl(U_k^L-a^2\bigr)>0,
\]
\[
\bigl(by-\tau_k^R\bigr)\bigl(V_k^L-\delta_a\bigr)
+
\bigl(ay+\tau_k^L\bigr)\bigl(v_k^R-b^2\bigr)>0.
\]
So the remaining local gap is no longer a question about the ratio \(r_k\); it is a cone-pairing
problem for the dual Weyl tails at the same block.

\smallskip
\noindent
\textup{(6)} Globally, the boundary slope is exactly
\[
\partial_y\lambda_*(x,y)=2s\sum_{j=1}^n p_j.
\]
At the same time, the weighted exact invariant gives
\[
s\sum_{j=1}^n \omega_j p_j
=
y\sum_{j=1}^n \omega_j |u_j|^2
>0.
\]
In the even case, all weighted suffix sums are also positive:
\[
s\sum_{j=\ell}^{2m}\omega_j p_j>0.
\]

\smallskip
\noindent
\textup{(7)} Therefore there are now two natural routes left:
either prove the local same-level tail pairings above, which would imply pointwise mirror-pair
positivity, or exploit the already-positive weighted mirror functional together with the local
current positivity to prove the weaker but sufficient inequality
\[
\sum_{j=1}^n p_j>0.
\]

\smallskip
\noindent
Thus the main unresolved issue has become genuinely local and finite-dimensional in the parity
Weyl variables, with a parallel global fallback through the weighted metric \(W\).
\end{remark}

\end{document}